\section{Discussion and Limitations}
\label{sec:discussion}

The empirical evaluation proves that while Mamba configurations can deliver high classification performance within tight microcontroller constraints—yielding $\sim$94\% accuracy on HAR and $\sim$89\% on KWS—their viability remains tightly coupled to specific compiler behaviors.

Our finding that model segmentation reduces peak working memory points toward an important architectural insight. We hypothesize that this memory reduction is an artifact of TFLM's internal arena allocation mechanism, which employs a greedy heuristic to schedule lifetime tensor offsets. By partitioning a single sequence graph into three smaller, isolated flatbuffer files, we effectively divide the tensor scheduling space into distinct subsets. This segmentation enables the greedy heuristic to achieve a more globally optimal memory layout than is possible when optimizing a massive, unrolled sequence graph. If this hypothesis holds, this memory-saving benefit will be highly sensitive to sequence topology and may not scale uniformly across different neural layer combinations.

From a structural standpoint, model segmentation serves as an essential enablement strategy rather than a pure performance optimization. Without this manual loop-partitioning scheme, deploying Mamba models for moderately long sequences is completely unfeasible under current TinyML compilers. This operational boundary was explicitly encountered during KWS compilation: the standard unrolled framework failed execution entirely due to internal TFLM operator constraints, which restrict the native \texttt{CONCATENATE} node to a maximum of 10 input dimensions, whereas the unrolled KWS architecture required merging 51 distinct temporal layers.

Several hardware and toolchain limitations qualify the findings of this study:
\begin{itemize}
  \item \textbf{Tooling and Compiler Immaturity:} Mainstream edge inference engines remain deeply optimized for standard feed-forward CNN layers. They lack native support for the recurrence mechanisms and dynamic matrix transformations unique to modern State-Space Models.
  \item \textbf{Dataset and Domain Boundaries:} This evaluation was limited to two representative temporal TinyML benchmarks (HAR and KWS). Consequently, these workloads do not capture performance trends in multidimensional domains, such as computer vision or high-frequency telemetry, where Mamba's linear scaling properties could yield different results.
  \item \textbf{Operator Bit-Width Constraints:} Our analysis was restricted to standard $FLOAT32$ and uniform $INT8$ quantization. Exploring sub-byte or mixed-precision bit-widths (e.g., $INT4$) represents a critical path to memory reduction, but such implementations are currently blocked by missing operator definitions in the core TFLM distribution.
  \item \textbf{Hardware and Ecosystem Vendor Lock:} Benchmarking was conducted exclusively on the Espressif ESP32 hardware family. Because microcontrollers vary dramatically in their custom vector extensions and non-volatile memory controller designs, evaluating these identical graphs on competing vendor architectures (e.g., ARM Cortex-M or STM32) would yield different latency and execution efficiency profiles.
\end{itemize}

Ultimately, manual sub-graph partitioning introduces significant system complexity, requiring host application code to be highly model-aware and to manually orchestrate internal state pointers. While this design layout successfully expands the upper bound of deployable model sizes, it introduces a severe latency penalty that complicates real-time applications. Resolving this trade-off requires the development of a unified, custom-quantized TFLM interpreter kernel that encapsulates the entire recurrent SSM loop natively. Such a specialized kernel would combine the memory efficiency of loop partitioning with the low overhead of an unrolled graph, and we leave its development to future work.
